% TOP_PROBLEM: {"id": 3000076, "title": "Scrambled Rota conjecture", "category_id": 2, "category": {"id": 2, "name": "combinatorics", "display_name": "Combinatorics", "description": "Counting problems, graph theory, discrete structures.", "slug": "combinatorics", "order_index": 2, "created_at": "2026-07-31T15:26:25.671Z"}, "status": "partially_solved", "problem_number": "AMR-029-0076", "set_id": 13}
% FIRST_POSTED: 2026-10-01
% ENTRY_KIND: statement_only
% UnsolvedMath ID 3000076; source catalogue code AMR-029-0076
% !TeX program = lualatex
% Definition and sources only; this file is not an LLM solution attempt.
\documentclass[11pt,a4paper]{article}
\usepackage[margin=25mm]{geometry}
\usepackage{fontspec}
\setmainfont{lmroman10-regular.otf}[BoldFont=lmroman10-bold.otf,
 ItalicFont=lmroman10-italic.otf,BoldItalicFont=lmroman10-bolditalic.otf]
\usepackage{amsmath,amssymb,mathtools}
\usepackage{unicode-math}
\setmathfont{latinmodern-math.otf}
\AtBeginDocument{\let\setminus\smallsetminus}
\usepackage{xurl,hyperref}
\hypersetup{colorlinks=true,urlcolor=blue}
\setcounter{secnumdepth}{0}
\setlength{\parindent}{0pt}
\setlength{\parskip}{5pt}
\setlength{\emergencystretch}{3em}
\begin{document}
\section*{Scrambled Rota conjecture}\label{3000076}
{\small UnsolvedMath ID 3000076 · AMR-029-0076 · Combinatorics · AMR Open Problem Lists}
\subsection{Definitions and mathematical statement}
Let $M=(S,\mathcal I)$ be a finite loopless matroid of rank $k\ge2$.
Independent sets form a hereditary family satisfying exchange;
bases are independent sets of size $k$. Looplessness means every
singleton is independent. Suppose the ground set can be partitioned
into $k$ bases, and hence has $k^2$ elements.

Now choose an arbitrary partition $S=A_1\sqcup\cdots\sqcup A_k$
with $|A_i|=k$. These new classes need not be independent.
The scrambled Rota conjecture asks for $k-1$ pairwise disjoint
bases $T_1,\ldots,T_{k-1}$ satisfying
\[
                              |T_j\cap A_i|=1
                                \quad(1\le i\le k,\ 1\le j\le k-1).
\]
Each requested basis is thus a full transversal of the new partition.

The original basis partition only certifies that the matroid has
enough disjoint bases; it need not have any alignment with the
new classes $A_i$. The $k-1$ transversals together use all but one
element of each class. The remaining $k$ elements are not required
to be a basis. Requesting $k$ disjoint transversal bases for arbitrary
new classes would be a stronger assertion than this conjecture.
No assumption of representability or additional independence of
the scrambled classes is included. The omitted case $k=1$ is
vacuous because it asks for zero transversals.
\subsection{Short English statement}
After arbitrarily regrouping the elements of a matroid that
splits into $k$ bases, must there still be $k-1$ disjoint transversal bases?
\subsection{Sources}
\begin{itemize}
\item[S1] ULAM.AI and the UnsolvedMath Contributors, supplied problems.json, record 3000076 (AMR-029-0076), AMR Open Problem Lists. Catalogue identity and source transcription; CC BY 4.0.
\url{https://huggingface.co/datasets/ulamai/UnsolvedMath}.
\item[S2] Egres Open - List of open problems, Open-problem page: Scrambled Rota conjecture. Original source indexed by AMR-029-0076.
\url{https://oldlemon.cs.elte.hu/egres/open/List_of_open_problems}.
\item[S3] EGRES Open, Scrambled Rota conjecture. Individual source page and explanatory remarks.
\url{https://lemon.cs.elte.hu/egres/open/Scrambled_Rota_conjecture}.
\end{itemize}
\end{document}
